\begin{center}
\textbf{Resumen}\par
\end{center}

\baselineskip

\noindent
\textbf{Título:} Diseño de un Aplicativo Instruccional de Realidad Aumentada Para la Asistencia en la Capacitación y Desarrollo de Actividades de Mantenimiento Preventivo\footnote[1]{Trabajo de Grado}

\noindent
\textbf{Autores:} Jesús David Galvis Carvajal y Juan Sebastián Carrillo Rangel\footnote[2]{Facultad de Ingenierías Fisicomecánicas. Escuela de Ingeniería Mecánica. Ingeniería Mecánica. \\ Director: German Orlando Romero Suárez Ph.D(C). en Ingeniería Mecánica. \\ Codirector: Diego Fernando Villegas Bermúdez Ph.D. en Ingeniería Mecánica.}

\noindent
\textbf{Palabras clave:} Realidad Aumentada, Mantenimiento Preventivo, Gestión de Información, Marcadores QR, ARSceneView, Guías Visuales, Animaciones 3D

\noindent
\textbf{Descripción:} En el presente proyecto se diseñó un aplicativo móvil de realidad aumentada orientado a la capacitación y al soporte en la ejecución de procedimientos de mantenimiento preventivo. El desarrollo inició con la identificación de requerimientos mediante una revisión bibliográfica y un análisis cualitativo de entrevistas realizadas a profesionales con experiencia en mantenimiento, cuyos resultados permitieron establecer las funcionalidades y características principales del sistema. Posteriormente, se desarrolló una aplicación para dispositivos Android empleando Kotlin, Jetpack Compose y la librería SceneView, incorporando módulos para la gestión de información técnica, consulta de documentación, programación de actividades, lectura de marcadores QR y visualización de modelos tridimensionales con guías secuenciales en realidad aumentada. La arquitectura del software fue estructurada con un modelo de gestión de estados basado en MVI, favoreciendo la organización y escalabilidad del sistema. Finalmente, se realizaron pruebas de funcionamiento y rendimiento utilizando un compresor de aire como caso de aplicación, evaluando aspectos relacionados con el inicio del aplicativo, la navegación entre pantallas y la correcta ejecución de las funcionalidades de realidad aumentada. Los resultados obtenidos evidencian el potencial de la realidad aumentada como herramienta de apoyo para la capacitación y la ejecución de actividades de mantenimiento preventivo.

\vspace{1em}

\clearpage

\begin{center}
\textbf{Abstract}\par
\end{center}

\baselineskip

\noindent
\textbf{Title:} Design of an Augmented Reality Instructional Application for Assisting Training and the Execution of Preventive Maintenance Activities\footnote[1]{Degree Work}

\noindent
\textbf{Authors:} Jesús David Galvis Carvajal and Juan Sebastián Carrillo Rangel\footnote[2]{Facultad de Ingenierías Fisicomecánicas. Escuela de Ingeniería Mecánica. Ingeniería Mecánica. \\ Director: German Orlando Romero Suárez PhD in Mechanical Engineering. \\ Co-Director: Diego Fernando Villegas Bermúdez PhD in Mechanical Engineering.}

\noindent
\textbf{Keywords:} Augmented Reality, Preventive Maintenance, Information Management, QR Markers, ARSceneView, Visual Guides, 3D Animations

\noindent
\textbf{Description:} This project involved the design of an augmented reality mobile application aimed at supporting training and the execution of preventive maintenance procedures. The development process began with the identification of system requirements through a literature review and a qualitative analysis of interviews conducted with professionals experienced in industrial maintenance. The findings served as the basis for defining the application's main functionalities and features. Subsequently, an Android application was developed using Kotlin, Jetpack Compose, and the SceneView library, incorporating modules for technical information management, documentation access, activity scheduling, QR marker recognition, and the visualization of three-dimensional models with sequential augmented reality guidance. The software architecture was structured following a Model–View–Intent (MVI) state management pattern, promoting organization and scalability. Finally, functionality and performance tests were conducted using an air compressor as the application case study, evaluating application startup, screen navigation, and the correct execution of the augmented reality features. The results demonstrate the potential of augmented reality as a tool for supporting both training and the execution of preventive maintenance activities.